\chapter{Stationäre Kosmologie}

\section{Einleitung}

Die WDBT+ führt zu einem stationären, nicht-expandierenden Universum. Die beobachtete Rotverschiebung ist keine Folge der Expansion des Raumes, sondern eine kumulative gravitative Wirkung auf Photonen während ihrer Reise durch das \( \Omega \)-Feld. Die DSTT-Drift (\( \dot{\beta} > 0 \)) bewirkt eine langsame Ausdehnung von Strukturen, die durch kontinuierliche Materieentstehung aus dem Quantenpotential kompensiert wird.

Dieses Kapitel beschreibt die Grundzüge dieser Kosmologie: Rotverschiebung, Galaxienrotation, Materieentstehung, Galaxienwachstum und die Interpretation des Sonnenwinds.

\section{Grundlagen der Rotverschiebung in der WDBT}

\subsection{Allgemeine Form}

Die allgemeine, aus den Weber-Kräften folgende Rotverschiebung für ein Lichtsignal, das eine radialsymmetrische Massenverteilung \( M(r) \) durchläuft, lautet:

\[
z = \frac{1}{c^2} \int_0^d \frac{G M(r)}{r^2} \, dr \qquad \text{(allgemeine Form)}
\]

Diese Gleichung gilt für jede Massenverteilung und enthält keine fraktale Dimension \( D \). Sie ist die Grundlage der Rotverschiebung in der WDBT.

\subsection{Rotverschiebung ohne Expansion}

In der WDBT+ ist die Rotverschiebung keine Doppler-Verschiebung durch Expansion, sondern eine kumulative gravitative Wirkung auf Photonen während ihrer Reise durch das \( \Omega \)-Feld. Die allgemeine Form (siehe oben) beschreibt genau diesen Effekt.

\subsection{Fraktale Dichteverteilung als Spezialfall}

Nimmt man speziell die fraktale Dichteverteilung der WDBT+ an (siehe Kapitel 6),

\[
\rho(r) = \rho_0 \left( \frac{r}{l_0} \right)^{D-3},
\]

so ergibt sich für die eingeschlossene Masse:

\[
M(r) = \int_0^r \rho(r') \cdot 4\pi r'^2 \, dr' = 4\pi \rho_0 l_0^{3-D} \frac{r^D}{D}.
\]

Einsetzen in die allgemeine Form liefert:

\[
z(d) = \frac{1}{c^2} \int_0^d \frac{G}{r^2} \cdot 4\pi \rho_0 l_0^{3-D} \frac{r^D}{D} \, dr
     = \frac{4\pi G \rho_0 l_0^{3-D}}{D c^2} \int_0^d r^{D-2} \, dr.
\]

Das Integral ergibt \( \int_0^d r^{D-2} dr = d^{D-1} / (D-1) \), und man erhält:

\[
z(d) = \frac{4\pi G \rho_0 l_0^{3-D}}{D(D-1) c^2} \, d^{\,D-1}.
\]

Dies ist die in der Literatur häufig zitierte Form. Sie ist ein Spezialfall der allgemeinen WDBT-Rotverschiebung für den Fall einer fraktalen Dichteverteilung. Für \( D \to 3 \) geht die fraktale Verteilung in eine konstante Dichte über.

Die effektive Hubble-Konstante wird skalenabhängig:

\[
H_{\text{eff}}(d) = c \cdot \frac{z(d)}{d} \propto d^{D-2}
\]

Mit \( D \approx 2,704 \) folgt \( H_{\text{eff}}(d) \propto d^{0,704} \). Die vollständige Herleitung findet sich in Anhang D.

Da \( \beta \) für Photonen von der Energie abhängt, zeigen verschiedene Frequenzen leicht unterschiedliche Rotverschiebungen. Dieser Effekt ist klein, aber mit Pulsar-Laufzeiten oder hochpräzisen Spektren messbar.

\section{Galaxienrotation ohne Dunkle Materie}

Flache Rotationskurven werden in der WDBT+ durch das \( \Omega \)-Feld und die anisotrope Trägheit der DSTT erklärt. Die radiale Dichteverteilung einer Galaxie ergibt sich aus der kontinuierlichen Materieentstehung und der DSTT-Drift.

\subsection{Dichteverteilung}

Die Materieentstehungsrate im fraktalen Raum ist (siehe Anhang J):

\[
\dot{\rho}_{\text{cre}}(r) \propto r^{D-3} = r^{-0,296}
\]

Im stationären Zustand führt dies über die Kontinuitätsgleichung zu einer radialen Dichteverteilung:

\[
\rho_{\text{Gal}}(r) \propto r^{-(3-D)} = r^{-0,296}
\]

\subsection{Rotationskurve}

Die kumulierte Masse innerhalb des Radius \( r \) ist:

\[
M(r) \propto \int_0^r \rho(r') r'^2 dr' \propto r^{2 + (3-D)} = r^{2,296}
\]

Daraus folgt die Rotationskurve:

\[
v(r) = \sqrt{\frac{GM(r)}{r}} \propto r^{0,648}
\]

Für \( D \approx 2,704 \) ergibt sich \( v(r) \propto r^{0,648} \) – dies ist eine fast flache Rotationskurve, konsistent mit Beobachtungen.

\section{Kontinuierliche Materieentstehung}

Die Materieentstehung wird durch die negative Vakuumenergiedichte des fraktalen Raumes vermittelt. Aus Anhang J folgt:

\[
\dot{\rho}_{\text{cre}}(r) = \kappa \cdot \frac{\hbar c}{l_0^4 c^2} \left( \frac{l_0}{r} \right)^{3-D} \propto r^{D-3}
\]

Die Materieentstehung ist also im Zentrum am stärksten und fällt mit einer gebrochenen Potenz ab, die direkt aus der fraktalen Dimension \( D \) folgt. Die vollständige Energiebilanz inklusive der negativen Vakuumenergiedichte findet sich in Anhang J.

\subsection{Drift und Einfang}

Die DSTT-Drift treibt die neu entstehende Materie nach außen. Die radiale Driftgeschwindigkeit wird in Anhang G hergeleitet:

\[
v_{\text{drift}}(r) = r\dot{\phi} \sqrt{1-\beta}
\]

Für Kreisbahnen als Näherung (\(\beta \approx 1\)) folgt die Skalierung \( v_{\text{drift}} \propto r^{-1/2} \).

Die Einfangwahrscheinlichkeit durch den Zentralkörper ist (Herleitung in Anhang G):

\[
\alpha_{\text{in}} = \left( \frac{r_c}{r_{\text{max}}} \right)^{3-D}
\]

Mit \( r_c \sim 200 \, \text{m} \) (Oberfläche des BEC-Objekts) und \( r_{\text{max}} \sim 5 \times 10^{11} \, \text{m} \) (äußerer Rand der Entstehungszone) ergibt sich \( \alpha_{\text{in}} \approx 0,001 \).

\section{Galaxienwachstum}

Während der Raum selbst nicht expandiert (das Universum ist stationär), wachsen gebundene Strukturen wie Galaxien durch kontinuierliche Materieentstehung und die DSTT-Drift. Aus der Drift-Gleichung \( dr/dt = v_{\text{drift}}(r) \) mit \( v_{\text{drift}} \propto r^{-1/2} \) folgt durch Integration:

\[
r(t) \propto t^{2/3}
\]

Die Materieentstehungsrate skaliert mit \( r^{3-D} \propto t^{2(3-D)/3} = t^{0,197} \). Daraus folgt für die kumulierte Masse der Galaxie:

\[
M_{\text{Gal}}(t) \propto t^{1,197}
\]

Der Exponent liegt nahe an 1, was ein nahezu lineares Wachstum der Galaxienmasse mit der Zeit bedeutet.

\subsection{Massenverhältnis Galaxie zu Kern}

Das Massenverhältnis von Galaxie zu zentralem BEC-Kern ist:

\[
\frac{M_{\text{Gal}}}{M_{\text{Kern}}} = \frac{1 - \alpha_{\text{in}}}{\alpha_{\text{in}}} \approx 1000
\]

Mit \( M_{\text{Kern}} = M_{\text{BEC}} \approx 3 \times 10^6 M_\odot \) ergibt sich \( M_{\text{Gal}} \approx 3 \times 10^9 M_\odot \) – typisch für große Galaxien.

\section{Die Sonne als Mikrokosmos}

Das Modell ist skalierbar. Für die Sonne (\( M = M_\odot \)) ergeben sich:

\[
r_{\text{max}} \sim 1,5 \times 10^{13} \, \text{m}, \quad r_c \sim 1,4 \times 10^{-11} \, \text{m}, \quad \alpha_{\text{in}} \approx 1,2 \times 10^{-7}
\]

\subsection{Interpretation des Sonnenwinds}

Im Standardmodell wird der Sonnenwind als Materieausstoß aus der Sonnenkorona gesehen. In der WDBT+ hingegen:

\begin{itemize}
    \item Materie entsteht in der Nähe der Sonne (bei \( r \sim 10^{-11} \, \text{m} \))
    \item Die DSTT-Drift treibt sie nach außen
    \item Nur ein winziger Bruchteil (\( \sim 10^{-7} \)) wird von der Sonne eingefangen
    \item Die Sonne verliert daher keine signifikante Masse
\end{itemize}

\section{Beobachtbare Signaturen: Der EHT-Schatten}

Ein BEC-Objekt mit \( M = M_{\text{BEC}} \) hat \( R = l_0 \approx 10^{-15} \, \text{m} \) – viel zu klein, um direkt abgebildet zu werden. Was das Event Horizon Telescope (EHT) zeigt, ist der Einflussbereich: die Region, in der Licht stark abgelenkt wird.

Der Photonenorbit liegt bei:

\[
r_{\text{ph}} \approx \frac{3GM}{c^2}
\]

Für \( M = M_{\text{BEC}} \):

\[
r_{\text{ph}} \approx 1,8 \times 10^{10} \, \text{m} \approx 0,12 \, \text{AU}
\]

Der Schatten ist also makroskopisch. Testbare Unterschiede zur ART:

\begin{itemize}
    \item Feinstruktur des Schattens
    \item Polarisationsmuster
    \item Zeitliche Variabilität (Flares)
\end{itemize}

\section{Vergleich mit dem Standardmodell}

Die folgende Tabelle fasst die Unterschiede zwischen der Standardkosmologie und der WDBT+ zusammen:

\[
\begin{array}{|l|l|l|}
\hline
\text{Phänomen} & \text{Standardmodell} & \text{WDBT+} \\
\hline
\text{Zentrale Objekte} & \text{Senken (verschlingen Materie)} & \text{Quellen (erzeugen Materie)} \\
\text{Materiefluss} & \text{Einwärts (Akkretion)} & \text{Auswärts (Drift)} \\
\text{Galaxien} & \text{Kollaps + Dunkle Materie} & \text{Wachstum von innen nach außen} \\
\text{Universum} & \text{Expandierend, Urknall} & \text{Stationär, ewig} \\
\text{Schwarze Löcher} & \text{Singularität + Horizont} & \text{BEC-Objekt, kein Horizont} \\
\text{Maximale Masse (Kern)} & \text{Keine natürliche Grenze} & M_{\text{BEC}} \approx 3 \times 10^6 M_\odot \\
\text{Dunkle Materie} & \text{Benötigt} & \text{Nicht benötigt} \\
\text{Dunkle Energie} & \text{Benötigt} & \text{Nicht benötigt} \\
\hline
\end{array}
\]

\section{Zusammenfassung}

Die WDBT+ führt zu einem stationären, nicht-expandierenden Universum:

\begin{itemize}
    \item Die Rotverschiebung ist keine Folge der Expansion, sondern kumulative Gravitation. Die allgemeine Form \( z = \frac{1}{c^2} \int \frac{GM(r)}{r^2} dr \) ist die Grundlage; die fraktale Dichteverteilung ist ein Spezialfall.
    \item Flache Rotationskurven werden ohne Dunkle Materie durch das \( \Omega \)-Feld und die anisotrope Trägheit der DSTT erklärt.
    \item Das Quantenpotential erzeugt kontinuierlich neue Materie, die durch DSTT-Drift nach außen getrieben wird.
    \item Galaxien wachsen von innen nach außen; das Massenverhältnis Galaxie/Kern folgt aus der fraktalen Raumdimension \( D \).
    \item Der Sonnenwind wird durch Materieentstehung in Sonnennähe erklärt – die Sonne verliert keine signifikante Masse.
    \item Der EHT-Schatten ist makroskopisch, aber die Feinstruktur unterscheidet sich von ART-Vorhersagen.
\end{itemize}

Die WDBT+ kommt ohne Urknall, Dunkle Materie und Dunkle Energie aus – und liefert ein konsistentes, singularitätenfreies Bild des Kosmos.